\chapter{The API Gateway}
\label{ch:gateway}

\section{The problem}

Without a gateway, the frontend must know six service addresses, each
service must implement its own authentication, and cross-cutting concerns
(CORS, rate limits, tracing headers) are done six times or not at all. A
gateway gives external clients \emph{one} address and centralizes the
edge concerns; behind it, services can trust their caller.

\section{Spring Cloud Gateway's model}

Three concepts compose every route:

\begin{description}
  \item[Route] a destination URI plus the conditions and rewrites below.
  \item[Predicates] conditions deciding whether a request matches:
  path, method, header, host. First matching route wins.
  \item[Filters] transformations applied on the way through:
  strip or add path segments, add headers, rate-limit, retry.
\end{description}

The gateway is built on the \emph{reactive} stack (Project Reactor +
Netty, \code{web-application-type: reactive} in its
\code{application.yml}) rather than servlet Tomcat. A proxy spends its
life waiting on other servers; an event-loop model holds thousands of
in-flight requests without one thread each. Operationally this is why
the gateway gets a smaller memory request than the MVC services in its
Deployment (384\,Mi vs.\ 512\,Mi) and boots faster --- which
Chapter~17 exploits with a shorter startup probe.

\section{The project's routes, annotated}

From \code{configurations/api-gateway.yml} (abridged --- the
course-service block repeats for students, assignments, enrollments,
submissions):

\begin{lstlisting}[language=yaml]
spring:
  cloud:
    gateway:
      routes:
        - id: auth-service
          uri: lb://auth-service          (*@\co{1}@*)
          predicates:
            - Path=/api/auth/**           (*@\co{2}@*)
          filters:
            - StripPrefix=2               (*@\co{3}@*)
        - id: auth-oauth2
          uri: lb://auth-service
          predicates:
            - Path=/oauth2/**             (*@\co{4}@*)
        - id: course-service-courses
          uri: lb://course-service
          predicates:
            - Path=/api/courses/**
          filters:
            - StripPrefix=2
app:
  jwt:
    secret: ${JWT_SECRET:mySecret...}     (*@\co{5}@*)
\end{lstlisting}

\begin{description}
  \item[\co{1}] \code{lb://} = ``resolve this name via the discovery
  client and load-balance.'' The name is the target's
  \code{spring.application.name}. Swap \code{lb://} for \code{http://}
  and you have a static proxy with no Eureka involved.
  \item[\co{2}] Path predicate. \code{**} matches any depth.
  \item[\co{3}] \code{StripPrefix=2} removes two segments:
  \code{/api/auth/login} arrives at auth-service as \code{/login}.
  This is why auth-service's SecurityConfig (Chapter~5) permits
  \code{/login}, not \code{/api/auth/login} --- the service is unaware of
  its public prefix.
  \item[\co{4}] No StripPrefix here: Google's OAuth2 redirect URI is
  registered with the full path, so it must arrive at auth-service
  unchanged.
  \item[\co{5}] The gateway holds the JWT \emph{validation} key. Same
  secret as auth-service, which \emph{signs} --- one shared secret, two
  roles, and in Kubernetes one Secret object mounted into both
  (Chapter~16).
\end{description}

\begin{pitfall}[404 from the service, 200 from curl-on-the-pod]
Symptom: calling \code{/api/courses/list} through the gateway returns
404, but the same endpoint answers when called directly on the service.
Cause: prefix mismatch --- either StripPrefix removed segments the
controller expects, or the controller mapping includes a prefix the
gateway already stripped. Debug by logging the \emph{rewritten} path:
enable \code{spring.cloud.gateway} debug logging and compare with the
controller's \code{@RequestMapping}.
\end{pitfall}

\section{Two routers: the gateway and the Ingress}

On Kubernetes this project has \emph{two} L7 routers in series: Traefik
(the Ingress controller, Chapter~18) at the cluster edge, then this
gateway. The division of labor is deliberate:

\begin{itemize}
  \item Traefik owns the \emph{host}: TLS later, one public entry,
  path-prefix fan-out. It sends \code{/api}, \code{/oauth2},
  \code{/login/oauth2} to the gateway and will send \code{/} to the
  frontend.
  \item The gateway owns the \emph{application edge}: Eureka resolution,
  JWT validation, per-route rewrites.
\end{itemize}

\begin{decision}[no side doors]
Milestone~3 had an Ingress sending \code{/courses} straight to
course-service --- useful scaffolding, but it bypassed JWT validation:
an unauthenticated path to protected data. It was deleted the day the
gateway was deployed. The rule it teaches: \emph{every} external route
must pass the component that enforces authentication, or the
authentication is decorative.
\end{decision}

\section{Advanced corners}

Worth knowing, not yet used here: \code{RequestRateLimiter} filter
(token bucket backed by Redis), \code{Retry} with backoff for idempotent
GETs, \code{CircuitBreaker} filter (Resilience4j) returning fallbacks
when a service is down, and default filters applied to all routes. Each
is a few YAML lines --- the gateway is where such policies belong,
because it is the one place every external request passes.

\section{Outside the project}

The same three-route idea in a Kubernetes-native gateway (no Eureka,
Service DNS instead) --- an nginx-ingress annotation style:

\begin{lstlisting}[language=yaml]
nginx.ingress.kubernetes.io/rewrite-target: /$2
# path: /api/courses(/|$)(.*)  ->  service: course-service
\end{lstlisting}

And the cloud-managed version: an AWS API Gateway or GCP API Gateway
route table doing predicates/filters as vendor config, with JWT
validation delegated to the provider (Cognito, IAP). The concepts
transfer 1:1; only the syntax and the billing change.

\begin{quiz}
\begin{enumerate}
  \item Follow \code{POST /api/auth/login} from browser to controller
  method on Kubernetes: name every hop and every path rewrite.
  \item Why must the OAuth2 routes not strip their prefix while the API
  routes must?
  \item The gateway validates JWTs; auth-service signs them. What
  breaks, and with what symptom, if the two ever hold different secrets?
  \item Why is a reactive stack the right choice for a proxy but not
  automatically for course-service?
\end{enumerate}
\end{quiz}
